\documentclass[10pt,a4paper]{amsart}
\usepackage[margin=19mm]{geometry}
\usepackage{amsmath,amssymb,mathtools}
\usepackage[scr=rsfs,cal=euler]{mathalfa}
\usepackage{xfrac}
\usepackage[unicode,hidelinks,bookmarks=false]{hyperref}
\newcommand{\ie}{i.e.\ }
\newcommand{\cf}{cf.\ }
\newcommand{\resp}{respectively\ }
\newcommand{\Subsection}{\subsection}
\providecommand{\nobreakdash}{\nobreak\hbox{-}}
\setlength{\emergencystretch}{2em}
\multlinegap0pt
\usepackage{amssymb}
\usepackage{mathtools}
\newtheorem{theorem}{Theorem}
\title{On Significant Pointwise Contractions and Annular Contractions}
\author{Ovidiu Popescu}
\date{Source: arXiv:2606.21319v1 (CC BY 4.0)}
\begin{document}
\maketitle

\noindent\emph{Source and licence.} On Significant Pointwise Contractions and Annular Contractions. Version arXiv:2606.21319v1 (CC BY 4.0), distributed by arXiv under the Creative Commons Attribution 4.0 International licence, http://creativecommons.org/licenses/by/4.0/, as recorded in the arXiv licence metadata for this exact version and shown on the abstract page https://arxiv.org/abs/2606.21319v1. Full licence notice: \texttt{licenses/arxiv-2606.21319-CC-BY-4.0.txt}.\par\smallskip

\noindent\textit{Editorial scope.} Counted unit: the article's theorem theorem (source0.tex lines 347--349). The article's complete original proof of it is delivered with it (source0.tex lines 351--378, one \texttt{proof} environment of 28 lines). No mathematics is written, repaired or re-derived: the statement and the proof are the article's own lines, in the article's own notation, and no retained line is substituted or re-flowed.  Retained uncounted as the source-local context the proof needs: lines 250--253; lines 329--332.  Nothing was omitted from any retained span, so the package carries no omission record.  No retained line contains a citation, so the wrapper carries no bibliography and no bibliography entry.  No retained line contains a figure environment, an image-inclusion command or drawing code, so no image is delivered and the package declares no asset dependency.  Editorial formatting only, in addition: an AMS \texttt{amsart} preamble that restates the article's own declarations and macros used by the retained lines, namely the environments \texttt{theorem} on separate counters, together with the standard packages those lines need.  The retained environments print in this document as Theorem 1 in order of appearance, whereas the article prints the same items with the numbering of its own full document.  The complete native source of the article as distributed in the arXiv e-print for this exact version is bundled unchanged as \texttt{sources/203265/source0.tex}.
	\begin{theorem}
		Let \((X,d)\) be a complete metric space and \(T : X \to X \) be an (SPC) mapping. Then, \(T\) has a fixed point.
		\label{Th2.25}
	\end{theorem}

		\begin{equation}
			\dfrac{1}{1+c}d(x,Tx) \leq d(x,y)\leq \dfrac{1}{1-c}d(x,Tx) \quad \text{implies} \quad d(Tx,Ty) \leq c d(x,y), 
			\label{***}
		\end{equation}

	\begin{theorem}
		Let \((X,d)\) be a complete metric space and \(T : X \to X \) be a (APC) mapping. Then, \(T\) has a fixed point.
	\end{theorem}

	\begin{proof}
		Let \(x_0 \in X\) be arbitrarily chosen and let \(\{x_n\}_{n=0}^{\infty}\) be the Picard iteration, i.e. the sequence defined by \(x_n = Tx_{n-1} = T^nx_0\). 
		
		Since \(\dfrac{d(x_{n-1},x_n)}{1+c} = d(x_{n-1}, Tx_{n-1}) \leq \dfrac{d(x_{n-1},x_n)}{1-c}\) for all \(n \geq 1\), taking \(x=x_{n-1}\) and \(y=x_n\) in \ref{***} we get 
		\[d(Tx_{n-1},Tx_n) \leq c d(x_{n+1},x_n).\]
		Hence, \(d(x_n,x_{n+1}) \leq c d(x_{n-1},x_n)\), for every \(n \geq 1\). Like in the proof of Theorem \ref{Th2.25}, we obtain that \(\{x_n\}\) is a Cauchy sequence. Since \(X\) is complete we have \(x_n \to x^* \in X\).
		
		If there exists \(n\geq 1\) such that 
		\[\dfrac1{1+c}d(x_{n-1},x_n) > d(x_{n-1},x^*) \quad \text{and} \quad \dfrac1{1+c}d(x_{n},x_{n+1}) > d(x_{n},x^*),\]
		then,
		\begin{align*}
			d(x_{n-1},x_n) &\leq d(x_{n-1},x^*)+d(x_{n},x^*) < \dfrac1{1+c}[d(x_{n-1},x_n)+d(x_{n},x_{n+1})]\\
			&\leq \dfrac1{1+c}[d(x_{n-1},x_n)+cd(x_{n-1},x_n)] = d(x_{n-1},x_n),
		\end{align*}
		which is a contradiction. Therefore, we have
		\[\dfrac{1}{1+c}d(x_{n-1},x_n) \leq d(x_{n-1},x^*) \quad\text{or} \quad \dfrac1{1+c}d(x_{n},x_{n+1})< d(x_{n},x^*).\]
		
		Since 
		\begin{align*}
			d(x_{n+p}, x_{n}) &\leq \sum_{i=0}^{p-1} d(x_{n+i+1},x_{n+i}) \leq \sum_{i=0}^{p-1} c^{i} d(x_{n+1},x_{n})\\
			& =(1+c+\dots+c^{p-1}) d(x_{n+1},x_{n}) < \dfrac{1}{1-c} d(x_{n+1},x_{n}),
		\end{align*}
		taking the limit as \(p \to \infty\), we get \(d(x^*,x_n) \leq \dfrac{1}{1-c} d(x_{n+1},x_n)\), for all \(n \geq 0\).
		
		Therefore, there exists a subsequence \(\{x_{n(k)}\}\) of \(\{x_n\}\) such that 
		\[\dfrac1{1+c}d(x_{n(k)},x_{n(k)}+1) \leq d(x_{n(k)},x^*) \leq \dfrac{1}{1-c}d(x_{n(k)},x_{n(k)}+1),\]
		for every \(k \geq 1\). Then \(d(Tx_{n(k)},Tx^*) \leq c d(x_{n(k)},x^*)\), by where we get, as \(k\to \infty\), that \(d(x^*, Tx^*) = 0\), i.e. \(x^*\) is a fixed point of \(T\).
	\end{proof}

\end{document}
